% !TEX root = ../main_zh.tex
\chapter{财政政策：李嘉图等价、拉弗曲线与政府债务}
\label{ch:fiscal}

凯恩斯的分析框架把一次财政扩张看作 IS 曲线的一次移动，至于这笔钱是靠税收还是靠借债来筹的，则交由某种未加说明的组合去解决。古典传统对此提出异议：筹资方式根本不是细枝末节。一个有前瞻性的家庭明白，今天的赤字就是明天的税收，并据此安排自己的计划——以至于在足够强的假设之下，征税与发债之间的选择对任何实际变量都不产生影响。本章在一个干净的两期模型里发展这一论断，它被称为\emph{李嘉图等价}（Ricardian equivalence）；进而辨明它所依赖的那些假设，由此看清它在现实中失效的种种理由；随后转向被这一等价结论重新框定、而非彻底了结的两个问题：税率究竟能走多高才会变得自拆台脚（拉弗曲线），以及政府债务是不是一种负担。

\section{李嘉图等价}

古典经济学家——斯密、李嘉图——认为市场能解决大多数问题，政府在很大程度上是不必要的。然而他们无法解释经济大萧条，凯恩斯正是在这里登场，提出用政府支出来撑起需求的主张。古典学派的回应——后来由新古典学派给出了严谨的微观基础——是：政府支出无论靠税收还是靠债务来筹措，都会挤出私人消费，因为理性的消费者看穿了赤字：今天卖出的一张债券，明天必须靠一笔税收来赎回。

考虑一个由单一代表性消费者和一个政府构成的经济体，存续两期。政府只在第一期支出，$G_1=G$，$G_2=0$，并用第一期税收 $t_1$ 与债券 $b$ 的某种组合来为之筹资，即 $G=t_1+b$。在第二期，它征收一笔税 $t_2=(1+r)b$ 来连本带息偿还债券。政府的预算可以用现值写成
\[
G = t_1 + \frac{t_2}{1+r},
\]
即税收的现值等于支出的现值。

\defn{李嘉图等价}{
\emph{李嘉图等价}是这样一个命题：对于给定的政府支出路径，在税收筹资与债券筹资之间的选择，不改变消费者的预算集——从而不改变消费与福利。重要的只是支出的现值，而非它如何筹措。}

消费者的偏好为 $U(C_1,C_2)=u(C_1)+\beta\,u(C_2)$，面临各期的预算约束
\[
C_1 = y_1 - t_1 - s, \qquad C_2 = y_2 - t_2 + (1+r)s,
\]
其中 $s$ 是私人储蓄。由于唯一的资产就是政府债券，在均衡中私人储蓄等于发行的债券，$s=b$。把两期的预算约束合并以消去 $s$，再代入政府预算约束，便得到消费者的一生预算约束
\[
C_1 + \frac{C_2}{1+r}
= y_1 + \frac{y_2}{1+r} - \left(t_1 + \frac{t_2}{1+r}\right)
= y_1 + \frac{y_2}{1+r} - G.
\]
消费的现值等于收入的现值减去支出的现值。$G$ 在 $t_1$ 与 $b$ 之间的拆分方式已经完全消失了。

\state{关键所在}{
消费者的一生预算只取决于 $G$ 本身，而与 $G$ 是靠税收还是靠债券筹得无关。两种 $G$ 相同的筹资方案给出相同的 $(C_1,C_2)$。这就是李嘉图等价。}

\ex[对数效用]{取 $u(C)=\log C$，于是消费者求解
\[
\max_{s}\; \log(y_1 - t_1 - s) + \beta\log\big(y_2 - t_2 + (1+r)s\big).
\]
之所以对 $s(=b)$ 求最优，是因为收入是外生给定的、税收由政府决定，消费者能够根据自己的偏好来确定储蓄量，而储蓄即投资于债券。一阶条件为
\[
\frac{1}{y_1 - t_1 - s} = \frac{\beta(1+r)}{y_2 - t_2 + (1+r)s},
\]
其解为
\[
b = s^{*} = \frac{\beta(1+r)(y_1 - t_1) - (y_2 - t_2)}{(1+\beta)(1+r)}.
\]
比较两种极端的筹资方案。在\emph{纯税收筹资}下，$G=t_1$，$b=0$，$t_2=0$，最优要求 $G = y_1 - \dfrac{y_2}{\beta(1+r)}$。在\emph{纯债券筹资}下，$b=G$，$t_1=0$，$t_2=(1+r)G$，代入后恰好给出同一个条件 $G = y_1 - \dfrac{y_2}{\beta(1+r)}$。两种情形下政府支出相同，因此 $C_1$ 与 $C_2$ 完全相同：等价成立；由同样的论证，它对介于两者之间的任何税收—债券组合也都成立。}

\subsection{那些假设，以及它们为何失效}

等价是一个定理，而像任何定理一样，它有自己的前提。

\asm{李嘉图等价的前提}{
\begin{enumerate}
\item \textbf{完全预见／理性：} 消费者预见到 $t_2=(1+r)b$，并对两期一并进行优化。
\item \textbf{消费者存活两期：} 从赤字中获益的，正是那个承担未来税收的同一主体。
\item \textbf{纳税人与债券持有人是同一个消费者。}
\item \textbf{税是总量税。}
\end{enumerate}}

每一条假设一旦失效，都会以一种富有启发的方式破坏等价。
\begin{itemize}
\item \emph{有限寿命。} 如果消费者活不到缴未来那笔税的时候，那么债券筹资的赤字会让当代人多消费，并迫使下一代少消费；消费不再在两代之间被平滑。（利他的遗赠通过把两代人联系起来，可以部分地恢复等价。）
\item \emph{纳税人 $\ne$ 债券持有人。} 这会产生两个效应。一是\emph{分配}效应：现在少买债券的人多消费，而将来不得不多买债券的人少消费。二是\emph{对储蓄与资本}的效应：当赤字未被额外的私人储蓄完全对冲时，政府借债吸走了本会用于投资的资金，使资本存量减少。
\item \emph{扭曲性税收。} 如果税不是总量税，而是收入相关的税或\emph{从价}税（例如累进税），它就会改变相对价格，把资源配置扭曲到偏离有效配置之外。现在征税制造的是当下的扭曲；先借债、将来再征税制造的是未来的扭曲。由于扭曲在何时发生现在变得要紧，筹资选择不再中性——这正是\emph{税收平滑}（tax smoothing）的依据：把扭曲性的税薄薄地摊到各期，并倚重债务来做到这一点。
\end{itemize}

\section{拉弗曲线}

一旦税收是扭曲性的，那么不仅它的时点，连它的\emph{水平}都变得要紧，提高税率也未必能提高收入。一个静态的劳动供给模型把这一点说清楚了。一位消费者在消费 $C$ 与闲暇 $l$ 之间权衡，效用为 $U(C,l)=C+\log l$，面临对劳动收入征收的比例税 $\tau$；设工资为 $w$、总时间单位化为 $1$，则预算约束为
\[
C = (1-\tau)(1-l)\,w.
\]
对 $l$ 最大化 $C+\log l$，一阶条件为 $-(1-\tau)w + 1/l = 0$，于是最优的闲暇与劳动供给为
\[
l^{*} = \frac{1}{(1-\tau)w}, \qquad 1 - l^{*} = 1 - \frac{1}{(1-\tau)w}.
\]
把预算约束改写为 $C + (1-\tau)w\,l = (1-\tau)w$ 便可看出：更高的 $\tau$ \emph{降低}了闲暇的有效价格 $(1-\tau)w$，于是由替代效应闲暇上升、劳动供给下降。税收收入是税率乘以（被征税的）劳动收入，
\[
R = \tau w\,(1-l^{*}) = \tau w - \frac{\tau}{1-\tau}.
\]
这两项有一个清晰的读法：$\tau w$ 是\emph{潜在收入效应}，即消费者一直工作时所能取得的收入；而 $-\dfrac{\tau}{1-\tau}$ 是\emph{扭曲效应}，其中 $\tau$ 是税率、$1-\tau$ 是它所诱发的扭曲。

\begin{figure}[h]
\centering
\begin{tikzpicture}
  \begin{axis}[
      width=12cm, height=7.2cm,
      scale only axis,
      axis lines=left,
      xmin=0, xmax=1,
      ymin=0, ymax=1,
      xtick={0,0.2,0.4,0.6,0.8,1},
      ytick={0,0.1,0.2,0.3,0.4,0.5,0.6,0.7,0.8,0.9,1},
      xlabel={Tax Rate $\tau$},
      ylabel={Labor Supply},
      x label style={font=\small},
      y label style={font=\small},
      tick label style={font=\footnotesize},
      axis line style={-Latex},
      clip=false,
    ]
    % Labor supply  LS = 1 - 1/((1-tau)*20),  w = 20, defined on [0,0.95]
    \addplot[red, line width=1.2pt, samples=200, domain=0:0.95]
      {1 - 1/((1-x)*20)};
    \node[red, anchor=south, font=\small] at (axis cs:0.32,0.92) {$1-l^{*}$};
  \end{axis}

  \begin{axis}[
      width=12cm, height=7.2cm,
      scale only axis,
      axis y line=right,
      axis x line=none,
      xmin=0, xmax=1,
      ymin=0, ymax=14,
      ytick={0,2,4,6,8,10,12,14},
      ylabel={Tax Revenue},
      y label style={font=\small},
      tick label style={font=\footnotesize},
      axis line style={-Latex},
      clip=false,
    ]
    % Revenue  R = 20*tau - tau/(1-tau),  defined on [0,0.95]
    \addplot[blue, line width=1.2pt, samples=200, domain=0:0.95]
      {20*x - x/(1-x)};
    \node[blue, anchor=south west, font=\small] at (axis cs:0.78,12.0) {$R$};
  \end{axis}
\end{tikzpicture}

\caption{拉弗曲线。随着税率上升，劳动供给（左轴）下降，先缓后陡；税收收入（右轴）画出一条倒 U 形，在税率效应占上风时上升，一旦不断收缩的税基占上风便转而下降。}
\label{fig:macro-30}
\end{figure}

求导可得 $\mathrm{d}R/\mathrm{d}\tau = w - 1/(1-\tau)^2$，它在使收入最大化的税率处为零；越过该点，收入便下降。图~\ref{fig:macro-30} 展示了这条倒 U 形：越过峰值之后，税基收缩得比税率上升得更快，于是收入下降。所谓\emph{禁止性}（prohibitive）税率，是指越过峰值的税率，此时劳动供给已经塌缩到提高税率反而降低总收入的地步；在那里，对最优配置的扭曲极为严重。

\state{现实启示}{
高税率有很强的扭曲性，因此政府应当把扭曲\emph{平滑}到各期，并更多地倚重债务。在战争时期，当支出必须骤增时，教科书式的处方是只适度提高税率，而用债券为其中的大部分融资——这恰恰是税收平滑的实践。}

\section{政府债务}

关于公共债务，有三个反复出现的问题。

\emph{政府债务可持续吗？}一个主权国家的债务不同于一个家庭的债务。只要国家存在，它的债务原则上就可以无限地滚动下去——发新债以赎旧债——而永远不必全额偿还。把债务与 GDP 相比是常见的做法，但并不完美：GDP 是流量而债务是存量；政府在债务背后持有大量资产；并且只要政权及其信用维系，债务就有担保。大体上说，债务是可控的，尽管即便财政尚有余地，评级机构的一次降级也可能袭来。

\emph{债务是经济的负担吗？}当政府借债来支出时，有些人今天少消费，而负担通过未来的税收在各代人之间分摊。通常的场合是经济萧条，政府支出以支撑需求。对于公众而言，持有这些债券本身就是一项正收益的资产，只要债务不断滚动，它就是安全的。

\emph{是否存在最优的债税比？}在李嘉图等价之下，税收与债券是等价的，因此并不存在最优的比例。但等价依赖于现实世界所违背的那些假设——主要是总量税——所以一旦税收是扭曲性的，税收平滑的论证便重新登场，于是确实存在一个真实的权衡。从历史上看，美国在第二次世界大战前后主要靠税收为自己筹资；自 1970 年代以来，天平已偏向债务，税收收入平稳，但开支扩张、赤字持续。

\section{社会保障}

社会保障支出是政府预算的一大部分，而它的设计回响着同一个筹资问题。

\defn{两种社会保障体制}{
\emph{确定给付制}（defined-benefit，即现收现付制）用今天劳动者的缴费来为今天的退休者发放——年轻人供养老年人。\emph{确定缴费制}（defined-contribution，即完全积累制）把每位劳动者的缴费拿去投资，并在其退休时连同投资收益返还给他——现在的自己供养未来的自己。}

关于现收现付制有一个关键事实：“账户”里的钱并不真的留作储备，它早已花在了当下的退休者身上。这样一种制度对投资风险是稳健的，但对人口结构的变化极为脆弱。随着人口老龄化——更少的劳动者供养更多的退休者，老年抚养比不断上升——现收现付制趋于赤字。补救之道无一令人愉快：提高缴费率会加重劳动者与企业的负担；削减给付在政治上不可能；推迟退休年龄或许在所难免。完全积累制避开了人口结构的挤压，却转而承担通胀与投资风险。
